\section*{अध्याय \ref{ch:g2:mixing-and-dissolving} --- मिलाना और घुलना}

\begin{solution}{exo:g2:mixing-and-dissolving:1}
चीनी और नमक \omterm{def:g2:mixing-and-dissolving:dissolve}{घुलते} हैं। रेत, खड़िया और कंकड़ नहीं \omterm{def:g2:mixing-and-dissolving:dissolve}{घुलते}।
\end{solution}

\begin{solution}{exo:g2:mixing-and-dissolving:2}
\omterm{def:g2:mixing-and-dissolving:dissolve}{विलयन} वह स्वच्छ द्रव है जो किसी ठोस के पानी में \omterm{def:g2:mixing-and-dissolving:dissolve}{घुलने} पर मिलता है।
उदाहरण: चीनी का पानी, नमकीन पानी (मीठी चाय भी एक \omterm{def:g2:mixing-and-dissolving:dissolve}{विलयन} है)।
\end{solution}

\begin{solution}{exo:g2:mixing-and-dissolving:3}
हाँ: जई, किशमिश, मेवे और बीज एक साथ रखे गए हैं। हर हिस्सा अब भी
दिखता है और बीना जा सकता है।
\end{solution}

\begin{solution}{exo:g2:mixing-and-dissolving:4}
नहीं। घुला हुआ ठोस पानी को स्वच्छ छोड़ता है और तली में कुछ नहीं रहता;
यहाँ पानी धुँधला है और एक परत नीचे बैठ गई है।
\end{solution}

\begin{solution}{exo:g2:mixing-and-dissolving:5}
वह स्वच्छ है (उसके आर-पार देखा जा सकता है), पर रंगहीन नहीं (वह गुलाबी है)। वह
\omterm{def:g2:mixing-and-dissolving:dissolve}{विलयन} है: \omterm{def:g2:mixing-and-dissolving:dissolve}{विलयन} का रंग हो सकता है।
\end{solution}

\begin{solution}{exo:g2:mixing-and-dissolving:6}
चाय का स्वाद मीठा है: चीनी अब भी उसमें है, इतने छोटे टुकड़ों में
पूरे पानी में फैली हुई कि दिख नहीं सकती।
\end{solution}

\begin{solution}{exo:g2:mixing-and-dissolving:7}
धूप पानी को सुखा देती है; समुद्र के पानी में घुला हुआ नमक पीछे
रह जाता है।
\end{solution}

\begin{solution}{exo:g2:mixing-and-dissolving:8}
``घुल गई।'' पिघलना वह है जो बर्फ़ का टुकड़ा अपने आप, बिना पानी के,
करता है; चीनी को पानी चाहिए, और वह अब भी चाय में है।
\end{solution}

\begin{solution}{exo:g2:mixing-and-dissolving:9}
रेत वाला बीकर धुँधला है और उसकी तली में परत है; चीनी का पानी स्वच्छ है
और उसकी तली में कुछ नहीं है।
\end{solution}

\begin{solution}{exo:g2:mixing-and-dissolving:10}
नल का पानी केवल पानी नहीं है: उसमें कुछ घुला हुआ है। पानी सूखने पर
वही ठोस सफ़ेद छल्ले के रूप में पीछे रह जाता है।
\end{solution}

\begin{solution}{exo:g2:mixing-and-dissolving:11}
वह हर गिलास की कुछ बूँदें एक साफ़, गहरे रंग की तश्तरी पर सूखने देती है।
शुद्ध पानी कुछ नहीं छोड़ता; नमकीन पानी सफ़ेद नमक छोड़ जाता है।
\end{solution}
